% GENERATED FILE. Edit canonical lessons, then run npm run generate:books.
% canonical-source-hash: fnv1a64:b4ee3f26b6c426c4
% canonical-lessons: MR-C200-snayu, MR-C200-jakham, MR-C200-vran, MR-C200-sardi, MR-C200-citthi

\chapter{Muscle, Wound, Scar, Head cold, Note}
\label{ch:mr-muscle-wound-scar-head-cold-note}
\hlchaptermodality{200}

\begin{chapteropening}
\textbf{By the end of this chapter:} \emph{I can say a muscle, a wound, a scar, a head cold and a note.}

Complete the last lesson of chapter 200: I can say a muscle, a wound, a scar, a head cold and a note.
\end{chapteropening}

\section[\texorpdfstring{\mr{स्नायू} (snāyū)}{स्नायू (snāyū)}]{\mr{स्नायू} (snāyū) — a muscle}

\label{lesson:MR-C200-snayu}



\begin{quote}
Before the new one: say the Marathi for to knit, then the Marathi for to harvest.
\end{quote}

\subsection*{You'll want to know: \mr{स्नायू}}
\textbf{\mr{स्नायू}} — \emph{snāyū} — \textquotedblleft{}a muscle\textquotedblright{}.

\begin{grammarlens}[title={What you've built}]
One more word for this chapter.
\end{grammarlens}

\subsection*{Your turn}
\begin{itemize}
\raggedright
  \item \emph{Say it:} \emph{snāyū}
  \item \emph{Say it:} \emph{snāyū}, once more
  \item \emph{Say it:} say \emph{snāyū}
  \item \emph{Recall:} say the Marathi for to knit, then the Marathi for to harvest, then say \emph{snāyū} again
\end{itemize}

\begin{tcolorbox}[breakable,colback=teal!4,colframe=teal!35!black,title={Before you move on}]
What does \mr{स्नायू} mean? (\textquotedblleft{}A muscle\textquotedblright{}.) Say it once more.
\end{tcolorbox}
\section[\texorpdfstring{\mr{जखम} (jakham)}{जखम (jakham)}]{\mr{जखम} (jakham) — a wound}

\label{lesson:MR-C200-jakham}



\begin{quote}
Before the new one: say the Marathi for to harvest, then the Marathi for a muscle.
\end{quote}

\subsection*{You'll want to know: \mr{जखम}}
\textbf{\mr{जखम}} — \emph{jakham} — \textquotedblleft{}a wound\textquotedblright{}.

\begin{grammarlens}[title={What you've built}]
One more word for this chapter.
\end{grammarlens}

\subsection*{Your turn}
\begin{itemize}
\raggedright
  \item \emph{Say it:} \emph{jakham}
  \item \emph{Say it:} \emph{jakham}, once more
  \item \emph{Say it:} say \emph{jakham}
  \item \emph{Recall:} say the Marathi for to harvest, then the Marathi for a muscle, then say \emph{jakham} again
\end{itemize}

\begin{tcolorbox}[breakable,colback=teal!4,colframe=teal!35!black,title={Before you move on}]
What does \mr{जखम} mean? (\textquotedblleft{}A wound\textquotedblright{}.) Say it once more.
\end{tcolorbox}
\section[\texorpdfstring{\mr{व्रण} (vraṇ)}{व्रण (vraṇ)}]{\mr{व्रण} (vraṇ) — a scar}

\label{lesson:MR-C200-vran}



\begin{quote}
Before the new one: say the Marathi for a muscle, then the Marathi for a wound.
\end{quote}

\subsection*{You'll want to know: \mr{व्रण}}
\textbf{\mr{व्रण}} — \emph{vraṇ} — \textquotedblleft{}a scar\textquotedblright{}.

\begin{grammarlens}[title={What you've built}]
One more word for this chapter.
\end{grammarlens}

\subsection*{Your turn}
\begin{itemize}
\raggedright
  \item \emph{Say it:} \emph{vraṇ}
  \item \emph{Say it:} \emph{vraṇ}, once more
  \item \emph{Say it:} say \emph{vraṇ}
  \item \emph{Recall:} say the Marathi for a muscle, then the Marathi for a wound, then say \emph{vraṇ} again
\end{itemize}

\begin{tcolorbox}[breakable,colback=teal!4,colframe=teal!35!black,title={Before you move on}]
What does \mr{व्रण} mean? (\textquotedblleft{}A scar\textquotedblright{}.) Say it once more.
\end{tcolorbox}
\section[\texorpdfstring{\mr{सर्दी} (sardī)}{सर्दी (sardī)}]{\mr{सर्दी} (sardī) — a head cold}

\label{lesson:MR-C200-sardi}



\begin{quote}
Before the new one: say the Marathi for a wound, then the Marathi for a scar.
\end{quote}

\subsection*{You'll want to know: \mr{सर्दी}}
\textbf{\mr{सर्दी}} — \emph{sardī} — \textquotedblleft{}a head cold\textquotedblright{}.

\begin{grammarlens}[title={What you've built}]
One more word for this chapter.
\end{grammarlens}

\subsection*{Your turn}
\begin{itemize}
\raggedright
  \item \emph{Say it:} \emph{sardī}
  \item \emph{Say it:} \emph{sardī}, once more
  \item \emph{Say it:} say \emph{sardī}
  \item \emph{Recall:} say the Marathi for a wound, then the Marathi for a scar, then say \emph{sardī} again
\end{itemize}

\begin{tcolorbox}[breakable,colback=teal!4,colframe=teal!35!black,title={Before you move on}]
What does \mr{सर्दी} mean? (\textquotedblleft{}A head cold\textquotedblright{}.) Say it once more.
\end{tcolorbox}
\section[\texorpdfstring{\mr{चिठ्ठी} (ciṭṭhī)}{चिठ्ठी (ciṭṭhī)}]{\mr{चिठ्ठी} (ciṭṭhī) — a note; a prescription}

\label{lesson:MR-C200-citthi}



\begin{quote}
Before the new one: say the Marathi for a scar, then the Marathi for a head cold.
\end{quote}

\subsection*{You'll want to know: \mr{चिठ्ठी}}
\textbf{\mr{चिठ्ठी}} — \emph{ciṭṭhī} — \textquotedblleft{}a note; a prescription\textquotedblright{}.

\begin{grammarlens}[title={What you've built}]
One more word for this chapter.
\end{grammarlens}

\subsection*{Your turn}
\begin{itemize}
\raggedright
  \item \emph{Say it:} \emph{ciṭṭhī}
  \item \emph{Say it:} \emph{ciṭṭhī}, once more
  \item \emph{Say it:} say \emph{ciṭṭhī}
  \item \emph{Recall:} say the Marathi for a scar, then the Marathi for a head cold, then say \emph{ciṭṭhī} again
\end{itemize}

\begin{tcolorbox}[breakable,colback=teal!4,colframe=teal!35!black,title={Before you move on}]
What does \mr{चिठ्ठी} mean? (\textquotedblleft{}A note; a prescription\textquotedblright{}.) Say it once more.
\end{tcolorbox}
